\section{Discussion and Limitations}
\label{sec:discussion}
\textbf{What the data show.} For the public V-JEPA ViT-L checkpoint we can state four things. First, its clean pairwise accuracy on IntPhys dev (81.7\%) is well below the published ViT-H figure (98\%), and it is not separable from a copy-last-frame pixel baseline (76.1\%; paired difference +4.7 pp, interval $-$1.1 to 10.6). Second, the verdict is unstable under a dither no viewer could see (21.5\% of pairs flip, 15.6\% strict), so any statement that a verdict is robust or fragile to an edit must be made against that floor. Third, our positive controls were not met: vertical flip raised surprise in 21.1\% of videos against a gate of 80\% (time reversal, not gating, 60.8\%), while a randomly initialised network did respond to vertical flip (87.8\%). Fourth, the pre-registered specificity rule returns ``detection-limited'': $A=$0.546 (0.501 to 0.591), and 0.520 (0.475 to 0.565) without global lighting, with a clean detection precondition of 0.722 against the required 0.85. Under the rule, specificity cannot be assessed here, and we make no claim in either direction: we do not claim that the signal tracks unpredictability rather than impossibility, nor the reverse.

\textbf{How these fit together.} The results are consistent with a reading in which the pairwise verdict on IntPhys dev relies largely on the discontinuities that the dataset's frame splices create. That the dataset is built from splices is a measured property of the dev videos (Section~\ref{sec:method}); that M1 relies on them is a hypothesis. A clip flipped vertically or reversed in time is internally consistent and has no such discontinuity. M1's surprise does not rise under vertical flip: in an exploratory paired check on the stored arrays its surprise decreases in 79.4\% of possible videos, while the random-weight network's increases in 87.8\%, and the cause is not established. Its accuracy on flipped pairs (83.9\%) does not fall. Time reversal, which also contains no splice, raises M1's surprise in 60.8\% of videos (lower confidence bound 51.7\%); this modest rise is not what a splice-only reading predicts, so the hypothesis is at best partial. A copy-last-frame baseline, which also responds to discontinuities, reaches a similar accuracy; M1 and B2 are not distinguishable at 180 pairs, which is not evidence that they are equivalent. We label the splice reading a hypothesis: our data do not rule out other explanations, and we did not test it directly. The random-weight network's response indicates that vertical flip changes something an untrained network also responds to, so the flip control is weaker than intended whatever M1 does.

\textbf{Exploratory splice analysis (post hoc, uncorrected, $n=180$).} Every one of the 180 impossible dev videos is a splice of its possible partners (120 single, 60 double). M1's accuracy is 91.2\% on the 137 pairs that the copy-last-frame baseline B2 gets right and 47.7\% on the 43 it gets wrong, and the per-pair margins of M1 and B2 correlate (Spearman 0.54, 0.42 to 0.65). M1 scores 75.0\% on single and 92.5\% on double splices, while B2 scores 75.0\% and 78.3\%. The window of peak surprise contains a switch frame in 73.9\% of impossible videos against 46.9\% expected by chance, and accuracy shows no monotone pattern across terciles of the splice position. Vertical flip changes few verdicts net (11 lost, 13 gained; McNemar $p=0.84$), whereas time reversal loses 23 pairs and gains 9 ($p=0.02$). These are associations. They do not separate ``M1 detects the cut'' from ``M1 and B2 both respond to post-splice scene differences'', and we do not take them as showing or establishing the splice reading. A causal test needs impossible videos that are not splices, which this dataset does not contain. On IntPhys~2 Main, M1 is at chance (47.6\%, 43.7 to 51.7); this agrees in direction with earlier reports of near-chance V-JEPA accuracy on that benchmark \cite{bordes2025,su2025}, and we draw no causal link.

\textbf{What this implies for how VoE results are reported.} A single pairwise accuracy cannot tell a model that tracks impossibility from one that tracks something correlated with it. We suggest reporting, next to accuracy, (i) the noise floor of the verdict under an imperceptible perturbation, (ii) the accuracy of trivial pixel baselines on the same pairs, and (iii) a positive control, with a randomly initialised model alongside it. None of these needs training. Our matching of abrupt edits on mean pixel difference was not sufficient: both pixel baselines landed below chance, so equal mean change does not mean equal visibility to every detector. A matched-event test with annotated event times and edits above the visibility threshold is a natural next step.

\textbf{Scope.} These conclusions concern one public checkpoint, on synthetic scenes, with frozen inference. V-JEPA~2 was run descriptively and is not H1-eligible, so nothing here extends to latent world models in general.

\subsection{Limitations}
\label{sec:limitations}
These limit every conclusion; details are in Section~\ref{sec:deviations}. (i) One public checkpoint, whose accuracy (81.7\%) falls short of the published ViT-H figure for an unexplained reason. (ii) The positive controls were not met, so we cannot show a response to violations beyond the dataset's frame splices; in IntPhys dev a ``violation'' is a splice. (iii) The specificity edits are near-threshold (median 0.33 of 255 grey levels) and placed at random windows for 178 of 180 pairs, and global lighting is a no-op in 84 of 180 entries. (iv) The H2 tie threshold ($1.09\times10^{-4}$) is small against the fp16 margin noise (up to $5.5\times10^{-4}$ per pair), so many effects of this size lie near the numerical floor. (v) Both pixel baselines are below chance, so $D_1$ and $D_2$ are artefact-prone. (vi) Invariance is descriptive: the dither floor exceeds 10\% and was measured in fp16 only, the zoom edit rests on 112 of 180 pairs, and the frame-rate edit is degenerate by our design error. (vii) M2 is H1-ineligible and descriptive, and M3 was not run. (viii) Seed variation covers nuisance parameters only and is not uncertainty; the scene bootstrap carries it. (ix) Synthetic scenes, one dataset for the main claims, frozen models, four abrupt-event types.
